% PROBLEM_FORMULATION.tex
% Canonical mathematical source for the BO problem and notation.
% Do not place long proofs here; those belong in THEORY.tex.

\section{Richly Conditioned Bayesian Optimization}
\label{sec:problem}

\paragraph{Purpose.}
This section defines the single problem solved by the paper.  The presentation is
intentionally Bayesian-optimization-first: structured conditioning is introduced
only after the standard sequential BO decision has been defined.

\subsection{Standard BO reference model}

At BO iteration $t$, let
\[
    \mathcal D_t = \{(x_i,y_i)\}_{i=1}^t
\]
denote the data collected so far, and let
\[
    P_{0,t,\theta}(df) = p_\theta(df\mid \mathcal D_t)
\]
be a tractable reference law for the unknown objective $f$, where $\theta$
collects the GP/kernel/likelihood hyperparameters.  We suppress $\theta$ and
write $P_{0,t}$ when it is fixed.  In the main paper
$P_{0,t}$ is a Gaussian-process posterior, although the theory only requires the
latent-Gaussian representation specified below when the structural influence
bound is instantiated.

\paragraph{Reference-model calibration scope.}
All subsequent conditioning and certification statements are conditional on
the supplied reference model $P_{0,t,\theta}$ at the current BO decision.  The
framework is agnostic to whether $\theta$ was pretrained or fixed, fitted from
ordinary BO observations, calibrated offline or periodically using the
structured factors, learned jointly with those factors, or obtained by an
approximate procedure.  In particular, genuine likelihood factors may inform
$\theta$ in a fully joint probabilistic model; learning or approximating
$\theta$ lies upstream of the present decision-conditioning guarantees.  If
factor-aware calibration itself requires expensive non-Gaussian inference, the
method does not remove that cost: its computational claim is only that it can
reduce the amount of rich-conditioning inference required to determine each BO
action, conditional on the calibrated reference model.  Calibration may in
practice be fixed, amortized, warm-started, or less frequent than acquisition
decisions, but those policies are not analyzed here.

The action domain is $\mathcal X$.  For a BO utility $u_x(f)$, the reference
acquisition is
\[
    \alpha_0(x) = \mathbb E_{P_{0,t}}[u_x(f)].
\]
The canonical example is expected improvement,
\[
    u_x(f) = (f(x)-y_t^\star)_+,
\]
where $y_t^\star$ is the current incumbent value.  The theory should allow
nonsmooth Lipschitz utilities such as EI and should not require $u_x$ to be
$C^1$ or $C^2$.

\subsection{Additional structured information}

Suppose BO also has access to structured information about $f$.  Represent this
information by conditioning factors
\[
    \{e_j(f)\}_{j=1}^N,
\]
with total energy
\[
    E_C(f) = \sum_{j=1}^N e_j(f).
\]
The fully conditioned law is
\begin{equation}
    \pi_C(df)
    = \frac{1}{Z_C}
      \exp[-E_C(f)]\,P_{0,t}(df),
    \label{eq:full-conditioned-law}
\end{equation}
where $0<Z_C<\infty$.

The factors may encode, for example, symmetry relations, monotonicity or shape
constraints, pairwise preferences, local physics/PDE residuals, conservation
relations, simulator-derived statements, or other structured information.
Because the factors may be nonlinear or non-Gaussian, $\pi_C$ is generally
non-Gaussian even when $P_{0,t}$ is a GP posterior.

The fully conditioned acquisition and BO action are
\begin{align}
    \alpha_C(x)
    &= \mathbb E_{\pi_C}[u_x(f)], \\
    x_C^\star
    &\in \arg\max_{x\in\mathcal X}\alpha_C(x).
\end{align}

The computational difficulty is that repeatedly evaluating and optimizing
$\alpha_C$ may require non-Gaussian inference in which all $N$ factors are
processed at every BO iteration.

\subsection{Decision-relevant active conditioning}

Let $S\subseteq[N]$ be an active subset of conditioning factors.  Define
\begin{align}
    E_S(f) &= \sum_{j\in S} e_j(f), \\
    H_S(f) &= \sum_{j\notin S} e_j(f),
\end{align}
and the active conditioned law
\begin{equation}
    \pi_S(df)
    = \frac{1}{Z_S}\exp[-E_S(f)]\,P_{0,t}(df).
    \label{eq:active-conditioned-law}
\end{equation}
Its acquisition is
\[
    \alpha_S(x)=\mathbb E_{\pi_S}[u_x(f)].
\]

The paper does \emph{not} require $\pi_S$ to be close to $\pi_C$ as a
distribution.  The target is the BO action.

Given a returned action $\widehat x$, define its full-conditioned acquisition
regret
\begin{equation}
    R_C(\widehat x)
    = \alpha_C(x_C^\star)-\alpha_C(\widehat x)
    = \sup_{x\in\mathcal X}
      \bigl\{\alpha_C(x)-\alpha_C(\widehat x)\bigr\}.
    \label{eq:full-regret}
\end{equation}
The computational goal is to return $\widehat x$ satisfying
\begin{equation}
    R_C(\widehat x)\le \epsilon
    \label{eq:target-certificate}
\end{equation}
while performing non-Gaussian inference with substantially fewer than $N$
conditioning factors whenever the structure of the problem permits it.

\subsection{Acquisition gaps are the canonical decision object}

For a challenger $x$ and proposed action $\widehat x$, define
\begin{align}
    F_{x,\widehat x}(f)
    &= u_x(f)-u_{\widehat x}(f), \\
    G_Q(x,\widehat x)
    &= \mathbb E_Q[F_{x,\widehat x}(f)]
     = \alpha_Q(x)-\alpha_Q(\widehat x),
\end{align}
for any law $Q$ over $f$.

Equation~\eqref{eq:full-regret} becomes
\[
    R_C(\widehat x)=\sup_{x\in\mathcal X}G_C(x,\widehat x).
\]
This pairwise gap is preferable to bounding acquisition values independently:
errors common to both actions can cancel, and only a challenger that can change
the ordering matters for the next BO decision.

\subsection{Active-to-full interpolation}

Introduce the path
\begin{equation}
    \pi_{S,s}(df)
    \propto
    \exp[-E_S(f)-sH_S(f)]\,P_{0,t}(df),
    \qquad s\in[0,1].
    \label{eq:interpolation}
\end{equation}
Then $\pi_{S,0}=\pi_S$ and $\pi_{S,1}=\pi_C$.
Under the regularity conditions stated in \texttt{THEORY.tex}, the acquisition
gap obeys the exact identity
\begin{equation}
    G_C(x,\widehat x)-G_S(x,\widehat x)
    =
    -\sum_{j\notin S}
    \int_0^1
    \operatorname{Cov}_{\pi_{S,s}}
    \!\left(F_{x,\widehat x},e_j\right)\,ds.
    \label{eq:gap-transport}
\end{equation}
This identity is the mathematical reason that conditioning relevance is
decision-specific.

\subsection{Screenable factor class}
\label{subsec:screenable}

The computational guarantee requires more structure than
Equation~\eqref{eq:gap-transport} alone.

Assume there is a finite Gaussian vector $Y$, defined jointly with $f$ under
the reference law, such that
\[
    Y\sim\gamma_t=\mathcal N(m_t,K_t),
    \qquad
    e_j(f)=\varepsilon_j(Y)
\]
for every conditioning factor.  Thus $Y$ is \emph{factor-sufficient}; it need
not determine the full function.  A deterministic representation
$f=\mathcal T_tY$ is a sufficient special case, but the reflection-symmetry
construction instead retains the reference conditional law of $f$ away from
the factor locations.

In the nonlinear-PDE instantiation, $Y$ is the complete discretized latent
field and the BO actions are field coordinates.  Rao--Blackwellization is
therefore the identity for the discretized EI observable, and the EI-gap
sensitivity for distinct actions $x$ and $\widehat x$ is supported only on
those two coordinates.  This is a special case of the factor-sufficient
formulation, not a restriction on the general framework.

For any integrable observable $g(f)$, define its Rao--Blackwellized observable
\[
    \bar g(Y)=\mathbb E_{P_{0,t}}[g(f)\mid Y].
\]
Let $\mu_{S,s}$ denote the $Y$-marginal along the active-to-full path:
\[
    \mu_{S,s}(dY)
    \propto
    \exp\!\left[-\sum_{j\in S}\varepsilon_j(Y)
                 -s\sum_{j\notin S}\varepsilon_j(Y)\right]\gamma_t(dY).
\]
Because the reweighting depends only on $Y$, the conditional law of $f$ given
$Y$ is unchanged from the reference law.  In particular,
\begin{equation}
    \operatorname{Cov}_{\pi_{S,s}}(g,e_j)
    =
    \operatorname{Cov}_{\mu_{S,s}}(\bar g,\varepsilon_j).
    \label{eq:rao-blackwell-covariance}
\end{equation}

Partition $Y$ into blocks $Y=(Y_1,\ldots,Y_B)$.  For a weakly differentiable
observable $\bar g$ of $Y$, let
\[
    L(\bar g)=(L_1(\bar g),\ldots,L_B(\bar g))^\top\in\mathbb R_+^B
\]
be a cheaply available block-sensitivity bound.  A factor family is called
\emph{screenable} for the BO utilities of interest if, for every active set and
every $s\in[0,1]$, there is an entrywise nonnegative influence operator $C_S$
such that
\begin{equation}
    \left|\operatorname{Cov}_{\mu_{S,s}}(\bar g,h)\right|
    \le L(\bar g)^\top C_S L(h)
    \label{eq:screenable-covariance}
\end{equation}
for the required Rao--Blackwellized gap and factor observables.  The notation
$C_S$ permits an active-set-dependent construction; a single conservative
$C_S\equiv C$ may instead be used when it is valid uniformly over $S$ and $s$.
The following computational requirements also apply:

\begin{enumerate}
    \item $C_S$ is obtained from model structure without sampling from the fully
          conditioned target;
    \item $L(e_j)$ can be computed from factor support, derivative bounds, or
          other inexpensive metadata without repeatedly evaluating omitted
          factors on posterior samples;
    \item application of $C_S$ to an aggregated sensitivity vector is cheaper
          than full conditioned inference.
\end{enumerate}

The abstract paper theory should be stated in terms of
Equation~\eqref{eq:screenable-covariance}.  A concrete block
Poincar\'e/Brascamp--Lieb/Dobrushin-style construction may be used to instantiate
$C_S$ for the experimental factor families, but that comparison theorem is
supporting machinery rather than the central contribution.

For expected improvement,
\[
    u_x(f)=(f(x)-y_t^\star)_+,
\]
the utility is 1-Lipschitz in $f(x)$ and weakly differentiable almost
everywhere.  Thus the gradient-based screenability argument should be stated
for weakly differentiable Lipschitz utilities rather than imposing unnecessary
smoothness.

\subsection{Structural, inference, and optimization error}

Define the structural omission error
\[
    \Delta_{\mathrm{struct}}(x,\widehat x)
    = |G_C(x,\widehat x)-G_S(x,\widehat x)|.
\]
Under screenability, use the bound
\begin{equation}
    B_{\mathrm{struct}}(x,\widehat x)
    =
    L(\bar F_{x,\widehat x})^\top C_S
    \sum_{j\notin S}L(e_j).
    \label{eq:structural-bound}
\end{equation}

Let $\widehat G_S(x,\widehat x)$ be an inference approximation satisfying,
on an event with specified confidence,
\begin{equation}
    |\widehat G_S(x,\widehat x)-G_S(x,\widehat x)|
    \le B_{\mathrm{infer}}(x,\widehat x).
    \label{eq:inference-interface}
\end{equation}
The confidence statement must be valid for the data-dependent leader,
challenger search, and adaptive refinement procedure; pointwise Monte Carlo
standard errors alone are not sufficient for a rigorous certificate.

Define the optimistic challenger envelope
\begin{equation}
    \Psi_S(x;\widehat x)
    =
    \widehat G_S(x,\widehat x)
    + B_{\mathrm{infer}}(x,\widehat x)
    + B_{\mathrm{struct}}(x,\widehat x).
    \label{eq:challenger-envelope}
\end{equation}
Suppose the challenger optimizer returns $x_c$ and a certified optimization
remainder $\eta_{\mathrm{opt}}$ satisfying
\[
    \sup_{x\in\mathcal X}\Psi_S(x;\widehat x)
    \le \Psi_S(x_c;\widehat x)+\eta_{\mathrm{opt}}.
\]
Then the stopping condition is
\begin{equation}
    \Psi_S(x_c;\widehat x)+\eta_{\mathrm{opt}}\le\epsilon.
    \label{eq:stop}
\end{equation}
The certificate theorem in \texttt{THEORY.tex} proves that
Equation~\eqref{eq:stop} implies Equation~\eqref{eq:target-certificate} on the
stated inference-confidence event.

A useful consequence is that the current leader need not itself be globally
certified.  A poor leader merely produces a stronger challenger and prevents
the stopping test from passing.  The global requirement belongs to the
optimistic challenger problem.

\subsection{Factor activation}

Let
\[
    h_U=\sum_{j\notin S}L(e_j).
\]
When $C_S=A_S^{-1}$, one structural solve
\[
    A_S w=h_U
\]
gives
\[
    B_{\mathrm{struct}}(x,\widehat x)=L(\bar F_{x,\widehat x})^\top w
\]
for all challengers.

For the worst challenger $x_c$, a factor-level upper-bound contribution can be
computed as
\[
    b_j
    = L(e_j)^\top C_S L(\bar F_{x_c,\widehat x}).
\]
If the certificate fails because structural error dominates, activate the
largest $b_j$ values (or $b_j/\mathrm{cost}_j$ when factors have heterogeneous
costs), recompute the active target, and repeat.  If inference or challenger
optimization error dominates, refine that component rather than adding factors.

\subsection{Continuous conditioning}

A continuous factor field may be written
\[
    E_C(f)=\int_{\mathcal Z} e(z,f)\,\mu(dz).
\]
The main paper need not develop this setting in detail.  The supplement may
replace the omitted sum in Equation~\eqref{eq:structural-bound} with the
componentwise integral
\[
    \int_{\mathcal Z} L(e_z)\,(\mu-\nu)(dz)
\]
for an active submeasure $\nu\le\mu$, or with the corresponding total-variation
bound for signed quadrature error.

\subsection{Scope and exclusions}

The exact perturbation identity is broad, but the cheap computational guarantee
is not.  The screenable class excludes, unless a separate valid influence bound
is supplied:

\begin{itemize}
    \item arbitrary black-box factors with no usable derivative, support, or
          coupling information;
    \item hard infinite-energy constraints for which the required covariance
          comparison is unavailable;
    \item factor interactions that destroy the needed Poincar\'e/comparison
          condition or make the influence operator dense and uninformative;
    \item models for which constructing $C_S$ is as expensive as full inference;
    \item discontinuous acquisition utilities without an alternative sensitivity
          argument.
\end{itemize}

The claim should therefore be stated as certified decision-relevant conditioning
for a structured, screenable class of richly conditioned GP/latent-Gaussian BO
models, not as a universal method for arbitrary conditioning statements.
